% Nouveau EGA V (cote 25): archive cote 25, site page 42 (PDF page 43).
% typescript. Draft compared: none.
% AI-transcribed by Claude Opus 5.5 (claude-opus-5-5, adaptive thinking, effort high), Message Batches API, batch ega25b (msgbatch_01A7n5YuXubUPNK1zsepi8RG), 2026-11-13; not yet checked by eye.
\scanpage{25/42}
\mspage{\hand{38}}

$\underline{Y}^{\mathrm{sing}} = P(\underline{\hand{\unsure{\nu}}}_{X/P}) \longrightarrow \hand{\unsure{\check{}}}\!P$ . Alors H'est ``orthogonal'' à x, et en parti-%
culier si dim H' = r-1 (ce qui équivaut par b) au fait que \xout{yx/x$\underline{Y}$xxxxx} \inserted[typed]{\unsure{supsi}}
y est point \inserted{géom.} singulier \uline{ordinaire} sur $Y_{\hand{\unsure{\xi|\zeta}}}$ ) x est uniquement déterminé
en fonction de H'.

Dém . Nous avons déjà prouvé a), et b) en est une conséquence
grâce aux développements du par.16. Reste à prouver c), qui est
encore conséquence d'une relation de compatibilité essentiellement
triviale. Je me borne ici à copier, dans le cas spécial qui nous
occupe (mais avec S quelconqque), le diagramme (D') du par.16, qui
prend la forme suivante (où nous écrivons $\underline{Y}^1$ pour $\underline{Y}^{\mathrm{sing}}$):

\diagram{
\begin{tikzcd}[column sep=small]
 & & 0 & & \\
 & & \underline{\Omega}^1_{X/S\,\underline{Y}^1} \arrow[u] & & \\
0 \arrow[r] & \underline{P}^1(\underline{O}(1,1))\otimes\underline{O}_{\underline{Y}^1} \arrow[r] \arrow[ur, "\hand{\unsure{\nu.\mu}}"] & (\underline{\Omega}^1_{X/S})_{\underline{Y}^1} + (\underline{\Omega}^1_{\check{P}/S})_{\underline{Y}^1} \arrow[r] \arrow[u] & \underline{\Omega}^1_{\underline{Y}^1/S} \arrow[r] & 0 \\
 & \underline{O}_{\underline{Y}^1}(-1,-1) \arrow[u] \arrow[r, "\hand{u}"] & \underline{\Omega}^1_{\check{P}/S}\ \underline{Y}^1 \arrow[u] \arrow[ur, "\hand{\nu}"'] & & \\
\underline{\omega}^{-2}_{X/Y}\otimes\underline{O}_{\underline{Y}^1}(-2,-2) \arrow[ur] & & 0 \arrow[u] & &
\end{tikzcd}
}

La relation d'orthogonalité c) ne fait qu'exprimer que le composé
$\over{V}{\nu}u$ dans ce diagramme est nul; je me dispense de développer plus
avant ce deuxième exercice de symtaxe. Je suppose que les gens avaient
l'habitude de prendre des coordonnées, et ils y voyaient peut-être plus
clair ....

\uline{Corollaire} 8.\over{\illegible{1}}{8}. Supposons S=Spec(k), \xout{Saxsaxtaxqusxx$\underline{Y}$xxxxx} \inserted{$f : X \to P$ une immersion} et X
irréductible \inserted{(\unsure{f propre} lisse sur k)}. \xout{xxxxx} Alors $\underline{Y}^{\mathrm{sing}}$ est lisse et irréductible; soit T

\margin{42}
